%  ============================================================================
% CHAPTER 2: MODEL AND THEORY
% ============================================================================

\section{Model and Theory}
As previously discussed, we can assume that agents can be modeled as particles in a closed system, where money is conserved.
Instead of exchanging energy, agents exchange money through random transactions. 
The system evolves through random transactions, forming a Markov process. 
The theoretical framework predicts that the wealth distribution converges to the Boltzmann-Gibbs distribution, with entropy increasing over time until equilibrium is reached.

\subsection{Conservation of Money}
One could assume that Boltzmann's derivation of the energy distribution can be applied to other statistical systems with a conserved quantity, such as money.
Indeed, we can consider the economy as a closed system with a fixed number of agents $\Agents$ and a total amount of money $\TotalMoney$.
Therefore, the total money in the system is conserved.
As a starting point, we can consider a simple model where agents exchange money through random transactions, and the total amount of money remains constant.

Introducing the concept of money conservation we can assume the model \textbf{without debt}, where agents cannot have negative balances.
So the conditions for this model are:
\begin{itemize}
    \item $\displaystyle \sum_{i=1}^{\Agents} m_i = \TotalMoney$.
    \item $m_i \geq 0$, for $i = 1, \ldots, \Agents$.
    \item $\delta m_i = -\delta m_j$, for $i \neq j$.
    \item $\delta_m < m_i$.
\end{itemize}

\subsection{Boltzmann-Gibbs Distribution of Money}
As argued in previous section, we can associate the money distribution
to the energy exchange in a physical system. 
For this reason, we can assume that the probability of finding an agent with money $m$ is given by the Boltzmann-Gibbs distribution:
\begin{equation}
    \Prob{m} = c e^{-m/\Temperature},
\end{equation}
where $c$ is a normalization constant and $\Temperature$ is the effective temperature of the system: $T = \TotalMoney / N$, which is average money per agent.

The \textbf{entropy} of the system is defined as:
\begin{equation}
    S = -\sum_{i=1}^{\Agents} \Prob{m_i} \log \Prob{m_i}.
\end{equation}
The \textbf{equilibrium} is reached when the entropy is maximized when the system "cools off" to the Boltzmann-Gibbs distribution.

\subsection{Transaction Mechanisms}
Many different kinds of transactions can be considered, such as:
\begin{itemize}
    \item \textbf{Constant transactions: $\Delta m = 1$} where a fixed amount of money is exchanged between agents. 
    \item \textbf{Fraction of Pair's Average:} A random fraction of the average money of the pair $\frac{m_i + m_j}{2}$ is exchanged:
                                                $\Delta m = r \cdot \frac{m_i + m_j}{2}$, 
    \item \textbf{Fractional transactions} where a fraction $\gamma$ of the agent's money is exchanged, the two agents will exchange
                                            only a portion of their money. This mechanism show how richer agents can exchange more money than poorer agents, and it is more realistic.
\end{itemize}

\subsection{The different models}
Different assumptions were made in the literature to model the money exchange between agents.
The aim it's to find a model that can reproduce the empirical wealth distribution observed in real economies.
\subsubsection{Model without debt}
The simplest model has constant transactions and no debt can be accumulated. This model leads to the Boltzmann-Gibbs distribution as previously discussed.
\subsubsection{Model with debt}
The exact mechanisms of limiting debt in the real economy are complicated and obscured.
In this scenario, debt is introduced in the model \cite{yakovenko2012statistical}. 




Therefore we can add a saving propensity $\lambda$ to the model, where agents save a fraction of their money before exchanging it.
